\chapter{Trazabilidad problema--requisito--solución}

Para que la oferta sea defendible ante el profesor y ante un cliente industrial, cada problema debe quedar conectado con un requisito, una función técnica y un indicador de verificación. La tabla siguiente evita que la propuesta parezca una lista de componentes sin justificación.

\begin{longtable}{P{0.20\textwidth}P{0.20\textwidth}P{0.26\textwidth}P{0.24\textwidth}}
\toprule
\textbf{Problema observado} & \textbf{Requisito derivado} & \textbf{Respuesta técnica} & \textbf{KPI de aceptación}\\
\midrule
\endhead
Operarios abandonan la estación para buscar útiles. & Reducir desplazamientos improductivos. & Misiones AMR de entrega y retorno de útiles usados. & Minutos de desplazamiento evitados por turno; misiones completadas sin intervención.\\
Kits o herramientas llegan tarde o incompletos. & Trazabilidad y confirmación de entrega. & Lectura RFID/QR, HMI de estación y registro de recepción. & Entregas a tiempo; incidencias de kit incompleto; entregas sin confirmación.\\
Riesgo FOD en zonas de trabajo. & Inspección preventiva documentada. & Cámara RGB-D, rondas FOD y evidencias fotográficas asociadas a misión. & Eventos FOD detectados; tiempo de respuesta; zonas inspeccionadas por turno.\\
Pasillos y estaciones con restricciones variables. & Navegación flexible y segura. & AMR omnidireccional con SLAM, LiDAR de navegación y escáner de seguridad independiente. & Paradas de seguridad; bloqueos de ruta; tiempo medio de replanificación.\\
Escalado de 40 a 80 HTP/mes. & Solución escalable multi-robot. & Gestor de flota, prioridades de misión y base de datos única. & Utilización de flota; disponibilidad; throughput logístico por turno.\\
Falta de evidencia fina para mejora continua. & Registro auditable de eventos. & Base de datos de misiones, trazabilidad y excepciones. & Porcentaje de misiones con registro completo; informes semanales generados.\\
\bottomrule
\end{longtable}

\section{KPIs recomendados para el piloto}

\begin{longtable}{P{0.28\textwidth}P{0.30\textwidth}P{0.30\textwidth}}
\toprule
\textbf{Indicador} & \textbf{Fórmula o criterio} & \textbf{Umbral de aceptación}\\
\midrule
\endhead
Éxito de misión & Misiones finalizadas / misiones lanzadas. & \(\geq 95\%\) en piloto controlado.\\
Disponibilidad AMR & Tiempo operativo / tiempo planificado. & \(\geq 90\%\) durante el piloto.\\
Entrega a tiempo & Entregas dentro de ventana / entregas totales. & \(\geq 90\%\) tras ajuste de rutas.\\
Registro completo & Misiones con kit, destino, hora y confirmación / total. & \(\geq 98\%\).\\
Incidentes de seguridad & Eventos con contacto o parada no controlada. & 0 incidentes con daño o contacto no previsto.\\
Aceptación de operarios & Encuesta breve de utilidad y carga operativa. & Media \(\geq 4/5\).\\
Tiempo improductivo recuperado & Horas estimadas antes--después por estación. & Tendencia positiva documentada, sin prometer plantilla menor.\\
\bottomrule
\end{longtable}

\section{Evidencia esperada}

La evidencia mínima para cerrar el piloto debe incluir mapa de rutas, listado de misiones, histórico de alarmas, registro de entregas RFID/QR, capturas FOD, encuesta de usuarios y comparativa de tiempos antes--después. Esta evidencia no sustituye una certificación aeronáutica, pero permite defender que la solución aporta valor medible.
